\newcommand{\p}{\mathcal{P}}
\chapterauthor{Joanna Balbus}
\chapter[Trwałość w układach \ldots]{Trwało\'s\'c w układach równa\'n reakcji -- dyfuzji z adwekcj\k{a}}
% \label{rffp} % Always give a unique label
% use \chaptermark{}
% to alter or adjust the chapter heading in the running head

\thispagestyle{empty}
\vspace{-8mm} {\large Joanna Balbus} \vspace{8mm}

	\section{Wprowadzenie}
	W ostatnich kilku latach wiele uwagi poświęca się modelom populacji w~środowisku jednorodnym lub niejednorodnym.
	Jednostajna trwałość, koegzystencja oraz wymieranie opisują istotne szczególne typy asymptotycznego zachowania rozwiązań powiązanych modeli. W niniejszej pracy badamy zagadnienie reakcji–dyfuzji typu Kołmogorowa z adwekcją
	
	\begin{equation}
		\label{1.1}
		% \tag{1.1}
		\begin{cases}
			u_{it}=\mu_i  u_{ixx}  -\alpha_i  u_{ix} + f_i(t,x,u_1,u_2)u_i, \quad t>0 \quad x\in (0,L),\quad i=1,2,\\
			u_i(0,x)=u_{i0}(x), \quad x\in (0,L),\quad i=1,2.
		\end{cases}
	\end{equation}
	z warunkami brzegowymi bez przepływu
	\begin{equation}
		\label{1.2}
		% \tag{1.2}
		\begin{cases}
			\mu_i u_{ix}(0,t) - \alpha_i u_i(0,t)=0, \quad t>0, \quad i=1,2, \\
			\mu_i u_{ix}(L,t) - \alpha_i u_i(L,t) = 0, \quad t>0, \quad i=1,2.
		\end{cases}
	\end{equation}
	
	W  literaturze istnieje wiele prac poświęconych równaniom reakcji–dyfuzji z adwekcją, zob. \cite{Hess,Huisman}, 
	\cite{Kolokolnikov,Lutscher},  \cite{Ma,Vasilyeva}, .
	Takie równania mogą być rozpatrywane z różnymi warunkami brzegowymi. W zależności od warunków brzegowych
	otrzymujemy odmienne scenariusze ekologiczne.
	
	Najbardziej ogólna postać warunków brzegowych to
	\begin{equation}
		\begin{cases}
			\mu_i u_{ix} - \alpha_i u_{i} =\tilde{b}_{i}\alpha_i u_i, \quad x=0,\quad i=1,2,\\
			\mu_i u_{ix} - \alpha_i u_{i} =-\hat{b}_{i}\alpha_i u_i, \quad x=L, \quad i=1,2.\\
		\end{cases}
	\end{equation}
	Parametry $\tilde{b}_i$ oraz $\hat{b}_i$ mierzą  tempo, z jakim osobniki opuszczają obszar odpowiednio na początku i końcu strumienia.
	
	Jeśli $\hat{b}_i=0$, $i=1,2$, otrzymujemy warunki brzegowe bez przepływu \eqref{1.2}, zob. \cite{Huisman} oraz \cite{Kolokolnikov}. Oznacza to, że żadne osobniki nie wchodzą do siedliska ani go nie opuszczają.
	Autorzy rozważali układ
	\begin{equation}
		\label{1.3}
		% \tag{1.3}
		\begin{cases}
			u_t=\mu_1u_{xx} -\alpha u_x+ (r(x) - u - v)u,\quad 0<x<L,\quad t>0,\\
			v_t=\mu_2  v_{xx} -\alpha v_x +(r(x)-u-v)v,\quad 0<x<L,\quad t>0,\\
			\mu_1 u_{x}(x,t) - \alpha u(x,t)=0, \quad t>0,\quad x=0,\quad x=L,\\
			\mu_2 v_{x}(x,t) - \alpha v(x,t) = 0, \quad t>0,\quad x=0,\quad x=L.
		\end{cases}
	\end{equation}
	
	Jeśli $\hat{b}=1$, mamy następujące warunki brzegowe
	
	\begin{equation}
		\label{1.4}
		% \tag{1.4}
		\begin{cases}
			\mu_1 u_x(0)  - \alpha u(0)=0,\\
			\mu_2 v_x(0) - \alpha v(0)=0,\\
			u_x(L)=0,\\
			v_x(L)=0.
		\end{cases}
	\end{equation}
	 Warunki brzegowe \eqref{1.4}
	nazywamy warunkami Danckwertsa.
	
	Zagadnienie \eqref{1.3} z warunkami \eqref{1.4} było rozważane w  \cite{Lutscher} i \cite{Vasilyeva} .
	
	Jeżeli $\hat{b} \to \infty$, otrzymujemy następujące warunki
	\begin{equation}
		\begin{cases}
			\mu_1 u_x(0)  - \alpha u(0)=0,\\
			\mu_2 v_x(0) - \alpha v(0)=0,\\
			u(L)=0,\\
			v(L)=0.
		\end{cases}
	\end{equation}
	
	Zagadnienie \eqref{1.3} z tymi warunkami brzegowymi było rozważane w \cite{Ma}.
	
	Dla $\alpha_i=0$, $i=1,2$, otrzymujemy następujące równanie reakcji–dyfuzji  
	\begin{equation}
		\label{1.5}
		% \tag{1.5}
		u_{it}=d_i u_{ixx} +f_i(t,x,u)u_i.
	\end{equation}
	
	W interpretacji biologicznej równanie \eqref{1.5} opisuje model populacji, w~którym tempo wzrostu uwzględnia jedynie losową dyspersję, zob. np.
	 \cite{Chafee,Hale,Rocha}.
	
	Zagadnienie \eqref{1.5} jest zwykle rozważane z warunkami Neumanna lub Dirichleta.
	
	Jednak na tempo wzrostu populacji wpływają różne czynniki zewnętrzne, takie jak dryf, wiatr, grawitacja.
	
	Model opisany równaniem \eqref{1.1} jest bardziej realistyczny niż modele opisane równaniem \eqref{1.2}.  
	Model \eqref{1.1} opisuje sytuację, w której osobniki podlegają jednokierunkowemu dryfowi, który wypycha je poza system, powodując w ten sposób spadek liczebności populacji. Doskonałymi przykładami środowisk z jednokierunkowym ruchem są ekosystemy wodne, takie jak strumienie czy rzeki, zob.  \cite{Lou,Speirs}. Osobniki zamieszkujące te systemy są znoszone w dół biegu rzeki przez ukierunkowany ruch. Pomimo tego znosu populacje opierają się wymywaniu i~utrzymują się przez wiele pokoleń w~takich środowiskach adwekcyjnych. To zjawisko biologiczne znane jest jako „paradoks dryfu” (zob. np. \cite{Anholt,Hershey,Huisman,Hsu,Muller1,Muller2,Pachepsky,Shigesada,Waters,Waters1}, ).
	
	Celem niniejszej pracy jest podanie warunków wystarczających, aby równanie \eqref{1.1} było trwałe przy pewnych podstawowych założeniach. W literaturze istnieją różne podejścia do dowodzenia globalnej atrakcyjności i trwałości w nieautonomicznych równaniach różniczkowych.
	Aby wykazać, że układ \eqref{1.1} jest trwały, korzystamy z metody pod- i nadrozwiązań.  Dynamika monotoniczna była rozważana w pracach  \cite{periodic,Hirsch,Yi,Smith}. W \cite{Zhao} Zhao rozważał układ monotoniczny i~podjednorodny prawie okresowy. 
	\vspace{2mm}
	
	Niniejsza praca jest zorganizowana następująco.
	W Rozdziale 2 wprowadzamy podstawowe założenia i lematy pomocnicze. Ponadto pokazujemy, że układ \eqref{1.1} ma jednoznaczne klasyczne rozwiązanie.
	Rozdziały 3 i 4 zawierają główne wyniki pracy. W Rozdziale 3 formułujemy średnie warunki trwałości układu \eqref{1.1}.
	
	\section{Preliminaries}
	
	W niniejszej pracy rozważamy konkurencyjny układ semiliniowych równań parabolicznych
	
	\begin{equation}
		\label{2.1}
		% \tag{2.1}
		\begin{cases}
			u_{it}=\mu_i  u_{ixx}  -\alpha_i  u_{ix} + f_i(t,x,u_1,u_2)u_i, \quad t>0,\; x\in (0,L),\quad i=1,2, \\begin{equation}2mm]
			\mu_i u_{ix}(0,t) - \alpha_i u_i(0,t)=0, \quad t>0, \quad i=1,2, \\begin{equation}1mm]
			\mu_i u_{ix}(L,t) - \alpha_i u_i(L,t) = 0, \quad t>0, \quad i=1,2,\\begin{equation}1mm]
			u_i(0,x)=u_{i0}(x), \quad x\in (0,L),\quad i=1,2,
		\end{cases}
	\end{equation}
	
	gdzie $u_i(t,x)$ oznacza gęstość $i$-tego gatunku w chwili $t$ i miejscu $x$ w~ograniczonym przedziale $[0,L]$,  
	$\mu_i>0$ są współczynnikami dyfuzji $i$-tego gatunku, a~$\alpha_i\ge 0$ oznaczają wielkość dryfu adwekcyjnego gatunku w~kierunku rosnącego $x$.  
	
	Z biologicznego punktu widzenia lewa krawędź (upstream) stanowi powierzchnię wody, a prawa (downstream) — dno.  
	Żadna z tych granic nie przepuszcza osobników, stąd narzucamy warunki zerowego strumienia.
	
	\vspace{3mm}
	
	Przyjmujemy następujące założenia:
	
	(A1) Funkcje $f_i:[0,\infty)\times [0,L]\times [0,\infty)^2 \to \mathbb{R}$ oraz ich pochodne pierwszego rzędu  
	$\frac{\partial f_i}{\partial t}$, $\frac{\partial f_i}{\partial u_j}$, $\frac{\partial f_i}{\partial x}$, $i=1,2$, są ciągłe.  
	Ponadto pochodne $\frac{\partial f_i}{\partial u_j}$, $i=1,2$, są ograniczone i jednostajnie ciągłe na zbiorze  
	$[0,\infty)\times [0,L]\times B$, gdzie $B$ jest dowolnym ograniczonym podzbiorem $[0,\infty)^2$.
	\newline
	Założenie (A1) jest standardowym warunkiem gwarantującym, że dla wystarczająco regularnej funkcji początkowej  
	$u_0(x)=(u_{0,1}(x),u_{0,2}(x))$, $x\in(0,L)$, gdzie $u_{0,1},u_{0,2}\ge 0$, układ \eqref{2.1} ma jednoznaczne, maksymalnie określone rozwiązanie  
	$u(t,x)=(u_1(t,x),u_2(t,x))$ dla $(t,x)\in [0,\tau_{\max})\times [0,L]$, $\tau_{\max}>0$, spełniające warunek początkowy.  
	Rozwiązanie jest klasyczne: wszystkie pochodne występujące w~równaniach i warunkach brzegowych istnieją oraz układ jest spełniony punktowo;  
	ponadto pochodne $\frac{\partial u_i}{\partial t}$ oraz $\frac{\partial^2 u_i}{\partial x^2}$ są ciągłe na $(0,\tau_{\max})\times [0,L]$ (por. Friedman).
	
	W konsekwencji silnej zasady maksimum, jeśli $u_{0,1}(x),u_{0,2}(x)\ge 0$ dla wszystkich $x\in [0,L]$ i nie są identycznie równe zero, to  
	$u_1(t,x),u_2(t,x)>0$ dla wszystkich $t\in (0,\tau_{\max})$ i $x\in (0,L)$.  
	Jeśli natomiast $u_{0,j}\equiv 0$ dla pewnego $j=1$ lub $2$, wówczas $u_j(t,x)=0$ dla wszystkich $t\in (0,\tau_{\max})$ i~$x\in (0,L)$.
	
	Przez rozwiązanie dodatnie układu \eqref{2.1} rozumiemy takie rozwiązanie $u(t,x)$, które spełnia $u_1(t,x),u_2(t,x)>0$ dla $t\in (0,\tau_{\max})$, $x\in (0,L)$.
	
	Regularność funkcji początkowych $u_0$ rozważa się zwykle w przestrzeniach Banacha ciągle zanurzonych w $C^1([0,L],\mathbb{R}^2)$  
	(see Henry \cite{Henry}). Ponieważ w pracy badamy zachowanie (dodatnich) rozwiązań dla dużych $t$, można — odpowiednio przesuwając czas początkowy —  
	założyć, że $u_{0,1},u_{0,2}$ są funkcjami klasy $C^2$ na $[0,L]$, dodatnimi w $(0,L)$ i spełniającymi warunki brzegowe.
	
	\vspace{5mm}
	
	Kolejne założenie jest standardowym założeniem ograniczoności.
	
	Dla $i=1,2$ funkcję $f_i(t,x,0,0)$ nazywamy wewnętrznym (intrinsic) tempem wzrostu $i$-tego gatunku.
	
	(A2) Funkcje $(t,x)\mapsto f_i(t,x,0,0)$ dla $i=1,2$ są ograniczone na $[0,\infty)\times [0,L]$.
	
	Niech
	\begin{equation}
	\overline{a}_i=\sup\{f_i(t,x,0,0): t\ge 0,\; x\in[0,L]\},
    \end{equation}
    \begin{equation}
	\underline{a}_i=\inf\{f_i(t,x,0,0): t\ge 0,\; x\in[0,L]\}.
	\end{equation}
	
	Dla ograniczonej ciągłej funkcji $c:[0,\infty)\to \mathbb{R}$ definiujemy średnią dolną:
	\begin{equation}
	m[c]=\liminf_{t-s\to \infty} \frac{1}{t-s}\int_{s}^{t} c(\tau)\,d\tau,
	\end{equation}
	oraz średnią górną:
	\begin{equation}
	M[c]=\limsup_{t-s\to \infty} \frac{1}{t-s}\int_{s}^{t} c(\tau)\,d\tau.
	\end{equation}
	
	Oznaczamy
	\begin{equation}
	m[f_i]=\liminf_{t-s\to \infty} \frac{1}{t-s}\int_{s}^{t} \min_{x\in[0,L]} f_i(\tau,x,0,0)\,d\tau,
	\end{equation}
	\begin{equation}
	M[f_i]=\limsup_{t-s\to \infty} \frac{1}{t-s}\int_{s}^{t} \min_{x\in[0,L]} f_i(\tau,x,0,0)\,d\tau.
	\end{equation}
	
	(A3) $m[f_i]>0$ dla $i=1,2$.
	
	Założenie (A3) jest dość szerokie: ograniczenia nałożone na wewnętrzne stopy wzrostu mają charakter całkowy —  
	istnieje $T>0$ takie, że
	\begin{equation}
	\int_{t}^{t+T} \min_{x\in [0,L]} f_i(\tau,x,0,0)\, d\tau >0
	\end{equation}
	dla każdego $t>0$.
	
	(A4) $\frac{\partial f_i}{\partial u_j}(t,x,u)\le 0$ dla wszystkich $t\ge 0$, $x\in [0,L]$, $u\in [0,\infty)^2$, $i\ne j$.
	
	Pochodna $\partial f_i/\partial u_j$ mierzy wpływ gatunku $j$ na tempo wzrostu gatunku $i$.  
	Układy \eqref{2.1} spełniające $\partial f_i/\partial u_j\le 0$ nazywane są konkurencyjnymi.
	
	(A5) Istnieje $\underline{b}_{ii}>0$ takie, że $\frac{\partial f_i}{\partial u_i}(t,x,u)\le -\underline{b}_{ii}$  
	dla wszystkich $t\ge 0$, $x\in [0,L]$, $u\in [0,\infty)^2$ oraz $i=1,2$.
	
	Z (A5) wynika, że mechanizm samoregulacji każdego gatunku jest na tyle silny, że poziomy populacji są ostatecznie ograniczone od góry przez wielkość zależną jedynie od parametrów układu.
	
	(A6) Istnieje $\overline{b}_{ij}\ge 0$ takie, że $\frac{\partial f_i}{\partial u_j}(t,x,u)\ge -\overline{b}_{ij}$  
	dla wszystkich $t\ge 0$, $x\in [0,L]$, $u\in [0,\infty)$ i $i,j=1,2$.
	
	\vspace{3mm}
	
	Niech $\varphi_i$, $i=1,2$, będzie jednoznacznie wyzanczoną funkcją własną zagadnienia
	
	\begin{equation}
		\label{2.6}
		% \tag{2.6}
		\begin{cases}
			\mu_i  \varphi_{ixx}  -\alpha_i \varphi_{ix}=\lambda_i \varphi_i, \quad x\in (0,L),\\
			\mu_i \varphi_{ix}(0) - \alpha_i \varphi_i(0)=0,\\
			\mu_i \varphi_{ix}(L) - \alpha_i \varphi_i(L)=0.
		\end{cases}
	\end{equation}
	
	
	
	
	Zagadnienie \eqref{2.6} ma dodatnią wartość własną, a odpowiadająca jej funkcja własna może być wybrana tak żeby była dodatnia na $[0,L]$ (por. \cite{Lutscher}).
	
	Definiujemy operator
	\begin{equation}
	Z[\psi_i]= \psi_{it}+\mu_i\psi_{ixx}-\alpha_i\psi_{ix}-f_i(t,x,u)\psi_i
	\end{equation}
	dla $x\in(0,L)$ oraz  
	$Z[\psi_i]=\mu_i\psi_{ix}-\alpha_i\psi_i$ na $x=0$ i $x=L$.
	
	Zakładamy, że dla $i=1,2$
	\begin{equation}
	\inf_{x\in [0,L]} \frac{u_{0,i}(x)}{\varphi_i(x)}\le 
	\sup_{x\in [0,L]} \frac{u_{0,i}(x)}{\varphi_i(x)}<\infty.
	\end{equation}
	
	Wprowadzamy rodzinę równań ODE, które będą przydatne przy badaniu dodatnich rozwiązań układu \eqref{2.1}.  
	Niech $u(t,x)$ będzie dodatnim rozwiązaniem \eqref{2.1} przy $f$ spełniającym (A1)–(A4).  
	Definiujemy $\xi_i(t)$, $t\ge 0$, jako dodatnie rozwiązanie problemu:
	
	\begin{equation}
		\label{2.7}
		% \tag{2.7}
		\begin{cases}
			\xi_i'(t)=\max_{x\in[0,L]} (f_i(t,x,0,0)-\lambda_i-\underline{b}_{ii}\xi_i)\xi_i,\\
			\xi_i(0)=\sup_{x\in[0,L]} u_i(0,x),\qquad i=1,2.
		\end{cases}
	\end{equation}
	
	\begin{lemma}
		Załóżmy (A1), (A2), (A4) oraz (A5).  
		Wówczas dla dowolnego dodatniego rozwiązania $u(t,x)$ układu \eqref{2.1} zachodzi
		\begin{equation}
		u_i(t,x)\le \xi_i(t)\varphi_i(x)\quad \text{dla }t\in[0,\tau_{\max}),\; x\in[0,L],\; i=1,2,
		\end{equation}
		gdzie $\xi_i(t)$ jest rozwiązaniem \eqref{2.7}.
	\end{lemma}
	
	\begin{proof}
	Wykażemy że  $\xi_i(t)$ jest  nadrozwiązaniem równania $u_i(t,x)$. 
	Z założenia (A5) 
	\begin{equation}
		f_i(t,x,u) \le \max_{x\in [0,L]} f(t,x,0,0)  - \underline{b}_{ii}u_i.  
	\end{equation}
	Zatem
	\begin{equation}
		\begin{split}
			Z[\xi_i \varphi_i] = \frac{\partial }{\partial t }(\xi_i \varphi_i) +\mu \frac{\partial^2 }{\partial x^2}(\xi_i \varphi_i) -
			\alpha_i \frac{\partial }{\partial x}\xi_i \varphi_i- f_i(t,x,u)\xi_i \varphi_i =\\
			\xi_i'(t) \varphi_i +\mu_i \xi_i \varphi_{ixx} - \alpha_i \xi_i \varphi_{ix}
			- f_i(t,x,u) \xi_i \varphi_i\\
			\geq (\max_{x\in [0,L]} f_i(t,x,0,0) - \underline{b}_{ii}\xi_i )\xi_i \varphi_i  +\mu_i \xi_i \varphi_{ixx} - \\
			\alpha_i \xi_i \varphi_{ix} - (\max_{x\in [0,L]} f_i(t,x,0,0) - \underline{b}_{ii}\xi_i -\lambda_i )\xi_i \varphi_i =0
		\end{split} 
	\end{equation}
	dla $i=1,2$ w przedziale $x\in(0,L)$.
	%
	%
	Ponadto
	\begin{equation}
		Z[\xi_i \varphi_i]= \mu_i \frac{\p(\xi_i \varphi_i)}{\p x} - \alpha_i\xi_i \varphi_i =0 
	\end{equation}
	w $x=0$ oraz $x=L$.
	Mamy także
	\begin{equation}
		\xi_i(0)\varphi_i(x)=\sup_{x\in[0,L]}\{\frac{u_i(0,x)}{\varphi_i(x)}\}\varphi_i(x)\le u_i(0,x).
	\end{equation}
	%
	Wtedy $u_i(t,x)\le \xi_i(t)\varphi_i(x)$ dla wszystkich $t\geq 0$ i $x\in [0,L]$.
	\end{proof}
	
	\begin{lemma}
		Załóżmy (A1)–(A4). Wówczas dla każdego maksymalnie określonego dodatniego rozwiązania $u(t,x)$ układu \eqref{2.1} zachodzi:
		
		(i) $\tau_{\max}=\infty$,
		
		(ii)
		\begin{equation}
		\limsup_{t\to\infty} u_i(t,x)\le \frac{\overline{a}_i}{\underline{b}_{ii}},\qquad i=1,2,
		\end{equation}
		przy czym zbieżność jest jednostajna względem $x\in[0,L]$.
	\end{lemma}
	
	\begin{proof}
		Ze standardowej zasady porównawczej dla równań różniczkowych zwyczajnych wynika że
		\begin{equation}
		\label{2.8}
		% \tag{2.8}
		\limsup_{t\to \infty} \xi_i(t)\le \frac{\overline{a}_i}{\underline{b}_{ii}}.
	\end{equation}
	Załóżmy że $\varphi_i$ jest znormalizowana tak że $\max_{[0,L]} \varphi_i(x)=1$, $i=1,2$. 
	Na podstawie Lematu 2.5 i z  \eqref{2.8} wynika że
	\begin{equation}
		\label{2.9}
		% \tag{2.9}
		\limsup_{t \to \infty}u_i(t,x) \le \limsup_{t \to \infty} \xi_i(t)\varphi_i(x) \le \frac{\overline{a}_i}{\underline{b}_{ii}}.
	\end{equation}
	Z Lematu 2.4 i \eqref{2.8} wynika że istnieje  $t_2\geq 0$ takie że
	 $u_i(t,x)$ jest ograniczona na odcinku $[t_1,\infty)$ przez 
	$\frac{\overline{a}_{i}}{\underline{b}_{ii}}+1$.
	Zatem rozwiązanie układu $\eqref{2.1}$ jest określone dla $t\in [0,\infty)$. To udowadnia  $(i)$. Dowód (ii) wynika z \eqref{2.9}.
	
	
	
	\end{proof}
	
	\begin{lemma}
		Niech $c:[t_0,\infty)\to\mathbb{R}$, $t_0>0$, będzie ograniczoną funkcją ciągłą, dla której istnieją liczby $\underline{c}>0$, $\overline{c}>0$ takie, że
		\begin{equation}
		-\underline{c}\le c(t)\le \overline{c}
		\end{equation}
		dla wszystkich $t\ge t_0$.  
		Niech $d>0$. Załóżmy ponadto, że istnieją $K>0$ i $\beta>0$ takie, że
		\begin{equation}
		\frac{1}{K}\int_{t}^{t+K} c(\tau)\,d\tau \ge \beta
		\end{equation}
		dla wszystkich $t\ge t_0$.  
		Wówczas dla dowolnego rozwiązania
		\begin{equation}
		\zeta'(t)=(c(t)-d\zeta)\zeta,\qquad \zeta(t_0)=\zeta_0>0,
		\end{equation}
		zachodzi
		\begin{equation}
		\frac{\beta}{d}e^{-K(\underline{c}+\beta)}
		\le \liminf_{t\to\infty} \zeta(t)
		\le \limsup_{t\to\infty} \zeta(t)
		\le \frac{\overline{c}}{d}.
		\end{equation}
	\end{lemma}
	
	\section{Średnie warunki trwałości}
	
	W niniejszej sekcji pokażemy, że przy pewnych warunkach dotyczących dolnych i~górnych średnich układ \eqref{2.1} jest trwały.
	%
	Najpierw wprowadzimy definicję trwałości dla układu \eqref{2.1}.
	\begin{definition}
		Układ \eqref{2.1} jest trwały (permanentny), jeśli istnieją dodatnie stałe $\delta_i$ takie, że dla każdego dodatniego rozwiązania
		$u(t,x) = (u_1(t,x),u_2(t,x))$ układu \eqref{2.1} zachodzi
		\begin{equation}
			\liminf_{t \to \infty} \frac{u_i(t,x)}{\varphi_i(x)} \geq \delta_i, \quad i=1,2,
		\end{equation}
		gdzie granica jest jednostajna względem $x \in [0,L]$, a $\varphi_i(x)$ jest funkcją własną zagadnienia \eqref{2.6}.
	\end{definition}
	%
	%
	\begin{theorem}
		Załóżmy (A1)–(A6). Jeśli
		\begin{equation}
			\label{3.1}
			% \tag{3.1}
			m[f_i] > \lambda_i + \frac{\bar{b}_{ij}(M[f_j] + \lambda_j)}{\underline{b}_{jj}}, \quad i=1,2,\ j\neq i,
		\end{equation}
		gdzie $\lambda_i$ jest główną wartością własną zagadnienia \eqref{2.6},  
		to układ \eqref{2.1} jest trwały.
	\end{theorem}
	
	\begin{proof}
		Ustalmy dodatnie rozwiązanie $u(t,x) = (u_1(t,x), u_2(t,x))$ układu \eqref{2.1}. Oznaczmy $\eta_i(t)$, $t \ge t_0$, jako dodatnie rozwiązanie zagadnienia początkowego
		\begin{equation}
			\label{3.2}
			% \tag{3.2}
			\begin{cases}
				\eta_i'(t) = \bigl(\min_{x\in [0,L]} f_i(t,x,0,0) - \lambda_i - \overline{b}_{ii}\eta_i - \overline{b}_{ij}\xi_j\bigr)\eta_i,\quad i\neq j,\\begin{equation}1mm]
				\eta(t_0)=\displaystyle\inf_{x\in[0,L]}\frac{u(t_0,x)}{\varphi(x)}.
			\end{cases}
		\end{equation}
		
		Pokażemy, że $\eta_i(t)\varphi_i(x)$ jest podrozwiązaniem dla $u_i(t,x)$.  
		Z założenia (A6) mamy
		\begin{equation}
			\begin{split}
				f_i(t,x,u) \ge \min_{x\in [0,L]} f_i(t,x,0,0)
				- \overline{b}_{ii} \eta_i(t)\varphi_i(x) - \overline{b}_{ij} u_j 
				\ge \\
				\min_{x\in [0,L]} f_i(t,x,0,0)
				- \overline{b}_{ii} \eta_i(t)\varphi_i(x)
				- \overline{b}_{ij} \xi_j.
			\end{split}
		\end{equation}
		
		Wobec tego
		\begin{equation}
			\begin{split}
				Z[\eta_i(t)\varphi_i(x)] 
				= \frac{\partial}{\partial t}(\eta_i(t)\varphi_i(x))
				+ \mu_i (\eta_i(t)\varphi_{ixx}(x)) 
				- \alpha_i (\eta_i(t)\varphi_{ix}(x)) \\
				\quad - f_i(t,x,u) (\eta_i(t)\varphi_i(x)) \\
				\le 
				(\min_{x\in [0,L]} f_i(t,x,0,0) - \lambda_i - \overline{b}_{ii}\eta_i - \overline{b}_{ij}\xi_j )\eta_i \varphi_i(x) \\
				\quad - \mu_i \eta_i(t)\varphi_{ixx}(x) + \alpha_i \eta_i(t)\varphi_{ix}(x) \\
				\quad - (\min_{x\in [0,L]} f_i(t,x,0,0)  - \overline{b}_{ii}\eta_i - \overline{b}_{ij}\xi_j) (\eta_i(t)\varphi_i(x)) \\
				\le -\lambda_i \eta_i \varphi_i - \eta_i(t) (\mu_i \varphi_{ixx} - \alpha_i \varphi_{ix}) = 0.
			\end{split}
		\end{equation}
		dla $t > t_0$, $x \in (0,L)$.
		
		Ponadto na brzegach
		\begin{equation}
			Z[\eta_i \varphi_i] = \mu_i \eta_i \varphi_{ix} - \alpha_i \eta_i \varphi_i = 0
		\end{equation}
		dla $x=0$ i $x=L$.
		
		A także
		\begin{equation}
			\eta_i(t_0)\varphi_i(x)
			= \inf_{x\in[0,L]}\frac{u_i(t_0,x)}{\varphi_i(x)}\varphi_i(x)
			\le u_i(t_0,x).
		\end{equation}
		
		Zatem
		\begin{equation}
			u_i(t,x) \ge \eta_i(t)\varphi_i(x), \qquad t \ge t_0,\ x\in [0,L].
		\end{equation}
		
		Wystarczy więc zastosować Lemma 2.6 do równania \eqref{3.2}, gdzie
		\begin{equation}
		c(t)= \min_{x\in [0,L]} f_i(t,x,0,0) - \lambda_i - \overline{b}_{ij}\xi_j,
		\qquad 
		d=\overline{b}_{ii}.
		\end{equation}
		
		Teraz pokażemy, że wielkości występujące w Lemma 2.6 można wybrać niezależnie od rozwiązania $u(t,x)$, przynajmniej dla dostatecznie dużych $t$.
		Łatwo zauważyć, że $c(t)$ ma górne ograniczenie niezależne od $u(t,x)$.
		
		Wybierzmy $\beta$ i $\beta'$ takie, że
		\begin{equation}
			\label{3.3}
			% \tag{3.3}
			0<\beta<\beta'< m[f_i] - \lambda_i + \frac{\bar{b}_{ij}(M[f_j]+\lambda_j)}{\underline{b}_{jj}}.
		\end{equation}
		
		Z \eqref{3.3} otrzymujemy
		\begin{equation}
		\frac{1}{L} \int_{t}^{t+L} \frac{\xi_j'(\tau)}{\xi_j(\tau)}\, d\tau
		= \frac{1}{L}\int_{t}^{t+L} \bigl(\max_{x\in[0,L]} f_j(\tau,x,0,0)
		+\lambda_j - \underline{b}_{jj}\xi_j(\tau)\bigr)\, d\tau.
		\end{equation}
		
		A więc
		\begin{equation}
			\label{3.4}
			% \tag{3.4}
			\underline{b}_{jj} M[\xi_j]
			= M[f_j] + \lambda_j
			- \frac{\ln \xi_j(t+L) - \ln \xi_j(t)}{L}.
		\end{equation}
		
		Z \eqref{2.9} wynika, że istnieje $L_0>0$ takie, że dla dowolnego dodatniego rozwiązania $u(t,x)$ istnieje $t_1\ge 0$ z własnością
		\begin{equation}
		\frac{|\ln \xi_j(t+L) - \ln \xi_j(t)|}{L}
		< \frac{(\beta'-\beta)\underline{b}_{jj}}{\max_{k\neq l}\overline{b}_{kl}}
		\end{equation}
		dla $j=1,2$, $t\ge t_1$, $L \ge L_0$.
		
		Wówczas
		\begin{equation}
			\begin{split}
				m\bigl[\min_{x\in [0,L]} f_i(t,x,0,0) 
				- \lambda_i - \overline{b}_{ij} M[\xi_j]\bigr] 
				= \\
				m[f_i] - \lambda_i 
				- \overline{b}_{ij} \left(
				-\frac{1}{\underline{b}_{jj}}
				\frac{|\ln \xi_j(t+L) - \ln \xi_j(t)|}{L}
				+ \frac{M[f_j]}{\underline{b}_{jj}} 
				+ \frac{\lambda_j}{\underline{b}_{jj}}
				\right)
				\ge \\
				m[f_i] - \lambda_i 
				- \frac{\overline{b}_{ij}(M[f_j] + \lambda_j)}{\underline{b}_{jj}}
				+ \frac{|\ln \xi_j(t+L) - \ln \xi_j(t)|}{L}
				> \beta,
			\end{split}
		\end{equation}
		dla dostatecznie dużych $t$.
		
		Pozostaje zatem zastosować Lemat 2.6, aby zakończyć dowód.
	\end{proof}
	
	\begin{thebibliography}{9999}
		%%
		
		\bibitem{Anholt} % drift paradox 
		B. R. Anholt,
		\emph{Density dependence resloves the stream drift paradox},
		Ecology 76 (1995) 2235 -- 2239
		%%
		%
		%
		%
		\bibitem{Chafee}
		N. Chafee,
		\emph{Asymptotic Behavior for Solutions of a
			One-Dimensional Parabohc Equation with
			Homogeneous Neumann Boundary Conditions},
		J. Differential Equations 18, 111 -- 134 (1975),
		%
		%%
		\bibitem{Hale}
		M. Chen, Xu-Yan Chen, J. K. Hale,
		\emph{Structural Stability for Time-Periodic
			One-Dimensional Parabolic Equations},
		J. Differential Equations 96, 355 -- 418 (1992),
		
		
		%
		%%
		%
		%%
		\bibitem{Henry}
		D. Henry,
		\emph{Geometric theory of semilinear parabolic equations},
		Lecture Notes in Mathemstics vol. 840 Springer Berlin New York (1981),
		%%
		\bibitem{Hershey} %drift paradox
		A. E. Hershey, J. Pastor, J. Peterson, W. G. Kling,
		\emph{Stable isotopes resolve the drift paradox for Baetis mayflies in an arctic river},
		Ecology 74,  2315 -- 2325 (1993),
		%%
		%%%
		%
		%%
		\bibitem{Hess}
		Hess, P,
		\emph{Spatial homogeneity of stable solutions of some periodic - parabolic problems with Neumann boundary conditions.},
		J. Differential Equations 68,  320 -- 331 (1987), 
		%%
		\bibitem{periodic}
		P. Hess, 
		\emph{Periodic - parabolic  boundary value problems and positivity},
		Pitman Research Notes in Mathematics Series 247 Longman Scientific and Technical New York 1991,
		%%
		%%
		%\bibitem{Hetzer}
		%G. Hetzer, W. Shen,
		%\emph{Uniform persistence, coexistence, and extinction in almost periodic / nonautonomous competition diffusion systems},
		%SIAM J. Math. Anal. 34(1),  204 -- 227 (2002),
		%%
		\bibitem{Hirsch} 
		M. W. Hirsch,
		\emph{Stability and convergence in strongly monotone dynamical systems},
		J. Reine. Angew. Math. 383,  1 -- 53 (1988).
		%%
		\bibitem{Huisman}  %drift paradox
		J. Huisman, M. Array\'as, U. Ebert and B. Sommeijer,
		\emph{How do sinking phytoplankton species manage to persist},
		Am. Nat. 159 245 -- 254,
		%%
		\bibitem{Hsu}  %drift paradox
		S. B. Hsu, Y. Lou,
		\emph{Single phytoplankton species growth with light and advection in a water column},
		SIAM J. Appl. Math. 70,  2942 -- 2974 (2010),
		%%
		\bibitem{Kolokolnikov}
		T. Kolokolnikov, C. Ou, Y. Yuan, Profiles of self-shading, sinking phytoplankton with finite depth, J. Math. Biol. 59, 
		105 -- 122 (2009),
		%%
		%%%
		%
		%
		%%
		\bibitem{Lou} %1-dimensional habitat
		Y. Lou, P. Zhou,
		\emph{Evolution of dispersal in advective homogeneous environment: The effect of boundary conditions.},
		J. Differential Equations 259,  141 -- 171, (2015)
		%
		\bibitem{Lutscher}
		F, Lutscher, Y. Lou,
		\emph{Evolution of dispersal in open advective environments},
		J. Math. Biol. 69 1319 -- 1342, 
		%
		
		
		\bibitem{Ma}
		F. Lutscher, E. McCauley, MA Lewis, \emph{Effects of heterogeneity on spread and persisence in rivers.}, 
		Bull. Math. Biol. 68(8), 2129 -- 2160,
		%
		
		%
		%
		%\bibitem{Mierczynski} %spectral theory
		%Mierczy\'nski, J., Shen, W.,
		%\emph{Spectral theory for random and nonautonomous parabolic equations and
		%	applications. In: Chapman  Hall/CRC Monographs and Surveys in Pure and Applied Mathematics.},
		%Chapman  Hall/CRC Boca Raton FL (2008),
		%%
		\bibitem{Muller1} 
		K. Muller,
		\emph{Investigations on the organic drift in North Swedish streams},
		Technical Report 34 Institute of Freshwater Research Drottningholm 1954,
		%%
		\bibitem{Muller2}
		K. Muller, 
		\emph{The colonization cycle of freshwater insects},
		Oecologica 53, 202 -- 207 (1982),
		%%
		%%
		\bibitem{Pachepsky}
		E. Pachepsky, F. Lutscher, R. Nisbet and M. A. Lewis, 
		\emph{Persistence, spread and the drift paradox},
		Theor. Pop. Biol. 67,  61 -- 73 (2005),
		%%
	
		%
		%
		%%
		\bibitem{Yi} %application in population dynamics
		W. Shen, Y. Yi,
		\emph{Convergence in almost periodic Fisher and Kolmogorov models},
		J. Math. Biol. 37 84 -- 102 (1998),
		%%
		%\bibitem{Shen} % z nabla
		%W. Shen, Y. Yi,
		%\emph{Almost automorphic and almost periodic dynamics in skew -- product semiflows.},
		%Mem. Amer. Math. Soc. 136 , no. 647 (1998),
		%%% 
		%
		\bibitem{Shigesada} %drift paradox
		N. Shigesada and A. Okubo,
		\emph{Analysis of the self -- shading effect on algal vertical distribution i natural waters},
		J. Math. Biol. 12, 311 -- 326  (1981),
		%%
		%
		%%
		\bibitem{Smith}
		H. Smith, 
		\emph{ Monotone dynamical system. An introduction to the theory of competitive and cooperative systems},
		Math. Surveys Monogr. vol 41 (1995) Amer. Math. Soc. Providence,
		%
		%%
		\bibitem{Speirs}
		D.C. Speirs, W.S.C. Gurney, 
		\emph{Population persistence in rivers and estuaries}, 
		Ecology 82, 1219 -– 1237 (2001),
		%%
		%
		%
		%
		\bibitem{Rocha}
		C. Rocha,
		\emph{Examples of Attractors
			in Scalar Reaction-Diffusion Equations} 
		J. Differential Equations 73, 178 -- 195 (1988),
		%%
		%%
		%
		%\bibitem{Wang} %[29] global attr
		%W. Shen, Y. Wang,
		%\emph{Carrying simpleces in nonautonomous and random competitive Kolmogorov systems},
		%J. Differential Equations 245,   1 -- 29  (2008),
		%%
		%%
		%\bibitem{Xiao}
		%Y. Wang, J. Xiao, D. Zhou,
		%\emph{Spatial - Homogeneity of stable solutions of almost - periodic parabolic equations with concave nonlinearity},
		%Proceedings of the American Mathematical Society April 2018 DOI: 10.1090/proc/14386,
		%%
		
		%%
		\bibitem{Vasilyeva}
		O. Vasilyeva, F. Lutscher, Population dynamics in rivers: analysis of steady states, Can. Appl. Math. Q. 18,
		439 -– 469 (2011),
		%%%
		
		%%
		\bibitem{Waters}
		R. F. Waters,
		\emph{The drift of stream insects},
		Ann. Rev. Entomol. 17,  253 -- 272 (1972),
		%%
		%
		\bibitem{Waters1} %drift paradox
		T. F. Waters,
		\emph{Seasonal patterns in production and drift Gammarus pseudolimnaeus in Valley Creek, Minnesota},
		Ecology 62,  1458 -- 1466 (1981),
		%
		%
		%
		%%
		\bibitem{Zhao}
		Zhao, X-Q, 
		\emph{Global attractivity in monotone and subhomogeneous almost periodic systems.},
		J. Differential Equations 187 , 494 -- 509 (2003).
		%
		%%
		%
		%
	\end{thebibliography}
	%
	